%% =====================================================================
%%  B3-T-03-01 -- what B3 contributes to "The Five Dangerous Marks"
%%  -------------------------------------------------------------------
%%  LAYER A: APPEND-ONLY. Written once, then left alone. Material found
%%  only here sits in the `unique` slot and is protected by rule R10.
%% =====================================================================
%% @kind: source
%% @id: B3-T-03-01
%% @book: B3
%% @topic: T-03-01
%% @locator: Law 19, pp. 137--147
%% @status: distilled
%% @unique: yes
%% @integrated: yes
%% @updated: 2026-09-19

\dossier{B3}{T-03-01}{\bookshort{B3}}{Law 19, pp. 137--147}

\begin{distilled}
  \item A preliminary typology of opponents, suckers and victims; distinguishing types is described as the highest form of the art, because it removes the need to coerce.
  \item The arrogant and proud man: touchy pride, vengeance out of all proportion to the slight, no sanity behind the over-reaction. Instruction is to flee on sensing it.
  \item The hopelessly insecure man: related but less violent and harder to spot; hurt simmers and arrives as small bites over a long period. Instruction is to disappear for a long time.
  \item Mr. Suspicion: sees the worst in everyone, genuinely unbalanced and therefore easy to deceive, but dangerous if his suspicion turns on you.
  \item The serpent with a long memory: no surface anger, calculates and waits, settles with cold shrewdness. Recognised by calculation across areas of life and by coldness. Instruction: crush completely or get out of sight.
  \item The plain, unassuming and unintelligent man: hardest to deceive, because taking a ruse requires imagination and a sense of possible reward. The cost is wasted time and resources rather than revenge. Test with a joke or story; a wholly literal reaction identifies the type.
\end{distilled}

\begin{unique}
  \item The full five-type typology, with the counter-intuitive finding that the apparently naive counterpart is the most expensive --- not dangerous, but unreachable, so the loss is time rather than injury.
  \item The joke-or-story test as a cheap discriminator for the literal type.
\end{unique}

\begin{terms}
  \item \emph{mark} --- the intended victim of a scheme.
  \item \emph{wolves from lambs} --- the book's shorthand for the whole task of typing a counterpart before engaging.
\end{terms}
